% !TEX root = userguide.tex

\def\headdate{28th January 2014}

\section{The eXtended Finite Element method (\xfem)}
\label{sec:xfem}

In this section some examples using the eXtended Finite Element method (\xfem)
will be described.
First some theoretical background will be given and then,
 for historical reasons, fluid flow examples (with
lots of details) and finally diffusion examples (with less
detail).

\subsection{Theoretical background}
\subsubsection{Basic ideas}
Partition of Unity method (PUM) \cite{Babuska97,Melenk06}\\
Generalized Finite Element method (GFEM)
 \cite{Strouboulis00a,Strouboulis00b,Strouboulis01}\\
The eXtended Finite Element method (\xfem) \cite{Belytschko99,Moes99}.\\
Fluid-structure interaction with \xfem\ \cite{Gerstenberger08}.\\
More to come \ldots

Both being based on the PUM, the difference between GFEM and \xfem\ is not very
clear.
Since \xfem\ has primarily been designed with discontinuities in mind and
our main application is discontinuities at material interfaces
we refer to 
methods based on the PUM using a finite element partition of unity as \xfem.

\subsubsection{A simple interpretation of \xfem\ for an enrichment
using Heaviside functions only}
\label{sec:xfem_interpretation}

Consider a body consisting of a single material with only smoothly varying
properties.
The field variables (like velocity, pressure in
case of a fluid) are not defined outside the domain of the body. Therefore,
we set all field variables to zero outside of the domain of the body. So,
considering both the inside and outside op the body,
the field variables show a jump across the boundary.

Assume we want to describe the motion of this body using a mesh
that is not boundary-fitted (see Fig.~\ref{fig:mesh_body}).
\begin{figure}[ht!]
   \begin{center}
   \includegraphics{mesh_body}
   \end{center}
   \caption{A material body and a regular (not boundary-fitted) mesh.}
    \label{fig:mesh_body}
\end{figure}
This means that the boundary is not aligned with
element boundaries.
We can numerically discretize this problem using \xfem\ with an enrichment
at the elements intersected by the boundary using a Heaviside function
(+1 on the material side and 0 on the outside). Elements fully outside the body
domain are ignored and nodes not attached to either
elements fully inside the body or to elements crossing the boundary are removed
from the system.

A simple interpretation of the \xfem\ procedure for this case is the following.
Consider, for example, the Poisson equation Eq.~\eqref{eq:poisson} defined on
a domain $\Omega$. The weak form is given by \eqref{eq:weakformpoisson}, where
the integration of the inner product Eq.\eqref{eq:innerproduct} is performed on
$\Omega$ only. In case of a boundary-fitted mesh the integration can be done
simply element by element. If we have the situation like in
Fig.~\ref{fig:mesh_body} elements fully outside $\Omega$ are ignored,
elements fully inside $\Omega$ are fully integrated as before, but elements
that are partially inside and outside need only to be integrated on the inside
part (see Fig.~\ref{fig:xfem_body}).
\begin{figure}[htbp!]
   \begin{center}
   \includegraphics{xfem_body}
   \end{center}
   \caption{A single material body for \xfem. Integration of the weak form is
only performed on the domain of the material body.}
    \label{fig:xfem_body}
\end{figure}
The degrees of freedom in the inside nodes have the same meaning as
before. The degrees of freedom in nodes outside of $\Omega$ but belonging
to elements intersected by the boundary are \emph{virtual degrees of freedom}.
They do not define the field variable at that point, but define the spatial
distribution of the field variable in the part of the element that is inside
the body. This is similar to virtual or ghost grid points in the finite
volume method.
Dirichlet boundary conditions need to be imposed differently
compared to the standard \fem. One possibility is to use a constraint
implemented using a Lagrange multiplier. Another possibility is to use
weak Dirichlet boundary conditions (see below).

If the outside is also filled with a material, nodes of the elements
intersected by the interface between the two materials will have degrees of
freedom both from the ``inside'' and ``outside'' (see
Fig.~\ref{fig:xfem_body2}). Furthermore, the two
materials need to be coupled. This can be done, for example,
by using constraints implemented with Lagrange multipliers.
Another possibility is to use weak coupling conditions.
\begin{figure}[htbp!]
   \begin{center}
   \includegraphics[scale=0.6]{xfem_body2}
   \end{center}
   \caption{Two material bodies for \xfem}
    \label{fig:xfem_body2}
\end{figure}

A crucial ingredient for the robustness and accuracy of the \xfem\ is the
integration on parts of an element domain.

An important difference with the fictitious domain method is that elements
or parts of elements not contributing to a body are removed from the system. This
means that in problems for fluid-solid interaction there is no fluid
inside the solid
and therefore also no false mass or viscous material response from a fictitious
fluid inside the solid.

\subsection{Weak Dirichlet boundary conditions for flow problems}
\label{sec:weakDflow}

Imposing Dirichlet conditions at the boundary of the material body (as in
Figure~\ref{fig:xfem_body}) is not straightforward, since the boundary is not
aligned with the edges/faces of the elements. One possibility is to use a
constraint implemented by Lagrange multipliers. Disadvantages of that
approach are:
\begin{itemize}
 \item The system has additional unknowns (Lagrange multipliers), which are
       located in ``nodes'' on the boundary and not in the nodes of the regular
       mesh.
 \item The required interpolation of the Lagrange multipliers is unclear and
       needs to satisfy the LBB condition for stability.
 \item The matrix structure is not very suited for fast iterative solvers.
\end{itemize}

An alternative approach is the use of weak Dirichlet boundary conditions
(see \cite{Choi11} for an overview of the possibilities and the method
described below).
The basic formulation is based on the open boundary formulation, where the open
boundary is applied to $\Gamma_\text{D}$, the Dirichlet boundary. For example, for
Stokes flow we find from Eq.~\eqref{eq:obcstokes}:
\begin{equation}
  \big(\ten D_v,2\eta\ten D\big)
    -(\nabla\cdot\vek v,p)
    -(\vek v,2\eta \ten D\cdot\vec n-p\vec n)_{\Gamma_{\text{D}}}=
(\vek v,\vek t_{\text{N}})_{\Gamma_{\text{N}}}+(\vek v,\vek f)
\qquad\text{for all }\vek v
\label{eq:embcstokes1}
\end{equation}
where $\vec n$ is the outwardly directed unit normal vector and 
the volume integrals evaluated on the material body only. This does not lead us
anywhere, because the Dirichlet condition:
\begin{equation}
    \vec u = \vec u_\text{D}\text{ on }\Gamma_\text{D}
\end{equation}
has not been imposed yet. We do that by adding to 
 Eq.~\eqref{eq:embcstokes1} a symmetric positive term:
\begin{multline}
  \big(\ten D_v,2\eta\ten D\big)
    -(\nabla\cdot\vek v,p)
    -(\vek v,2\eta \ten D\cdot\vec n-p\vec n)_{\Gamma_{\text{D}}}\\
 +\kappa M_{km}^{-1}(\vec v\vec n,\phi_k)_{\Gamma_\text{D}}:(\phi_m,[\vek n(\vec u-\vec
u_\text{D})+(\vec u-\vec u_\text{D})\vec n])_{\Gamma_\text{D}}\\
=(\vek v,\vek t_{\text{N}})_{\Gamma_{\text{N}}}+(\vek v,\vek f)
\qquad\text{for all }\vek v
\label{eq:embcstokes2}
\end{multline}
where $\phi_k$ is a shape function defined on all elements and
$M_{km}=(\phi_k,\phi_m)$.
We have used summation convection for dummy indices $k$ and $m$,
both running over all nodes.
 The parameter $\kappa>0$ has a dimension
identical to a viscosity and can be chosen freely, but a value near the actual
viscosity gives the best results.

Remarks:
\begin{enumerate}
   \item The full weak form is obtained by including the term (continuity eq.)
\begin{equation}
   \ldots -(q,\nabla\cdot\vek u) \dots
\end{equation}
to the left-hand side of the weak form.
   \item The method is an extension of the Gerstenberger-Wall scheme
        \cite{Gerstenberger10} for use with non-Newtonian fluids
        \cite{Choi11}. The original Gerstenberger-Wall formulation uses
        an additional field that is numerically eliminated when building the
        system. In the formulation above, the field has been eliminated from the
        formulation itself to obtain the so-called primary formulation.
   \item For Stokes $\kappa=\eta$ is optimal.
   \item The interpolation order of the shape function $\phi_k$ needs to be at
least equal to the velocity interpolation and in the implementation we
 take identical shape functions.
   \item The inverse of the mass matrix $M_{km}$ needs to be computed. Since we
will use discontinuous shape functions $\phi_k$, this can be done at element
level.
   \item Although $\phi_k$ is defined in all elements, the contribution of the
additional term is only non-zero for elements that are crossed by the boundary.
   \item The method is consistent.
   \item The method is not limited to the Stokes problem and can be applied to
         other flow problems (Navier-Stokes, viscoelastic, \dots)
         by adding an identical additional term to the open
         boundary formulation for the problem at hand.
   \item The method is non-symmetric even for a symmetric problem, such as the
         Stokes problem. 
This can be remedied
easily by adding the transposed boundary terms to the
weak momentum balance \cite{Choi11}:
\begin{multline}
  \big(\ten D_v,2\eta\ten D\big)
    -(\nabla\cdot\vek v,p)
    -(q,\nabla\cdot\vek u)
    -(\vek v,2\eta \ten D\cdot\vec n-p\vec n)_{\Gamma_{\text{D}}}
-(2\eta \vec n\cdot\ten D_v-q\vec n,\vec u-\vec u_\text{D})_{\Gamma_{\text{D}}}\\
 +\kappa M_{km}^{-1}(\vec v\vec n,\phi_k)_{\Gamma_\text{D}}:(\phi_m,[\vek n(\vec u-\vec
u_\text{D})+(\vec u-\vec u_\text{D})\vec n])_{\Gamma_\text{D}}\\
=(\vek v,\vek t_{\text{N}})_{\Gamma_{\text{N}}}+(\vek v,\vek f)
\qquad\text{for all }\vek v
\label{eq:embcstokes4}
\end{multline}
Note, that $\ten D_v=(\nabla\vek v+\nabla\vek v^T)/2$.
Contrary to what would
be expected, the non-symmetric method is more stable than the symmetric method
for lower values of $\kappa$ \cite{Choi11}.
   \item The positive term stabilizes the method. Another possibility for
stabilization is the Baumann-Oden approach, which adds a \emph{negative}
         transposed open
         boundary term:
\begin{multline}
  \big(\ten D_v,2\eta\ten D\big)
    -(\nabla\cdot\vek v,p)
    -(\vek v,2\eta \ten D\cdot\vec n-p\vec n)_{\Gamma_{\text{D}}}
+(2\eta \vec n\cdot\ten D_v-q\vec n,\vec u-\vec u_\text{D})_{\Gamma_{\text{D}}}\\
=(\vek v,\vek t_{\text{N}})_{\Gamma_{\text{N}}}+(\vek v,\vek f)
\qquad\text{for all }\vek v
\label{eq:embcstokes6}
\end{multline}
where the pressure part might be left out, since the method is non-symmetric
anyway.
Note the antisymmetry of the boundary terms due to the sign change of the term
with the velocity ``jump'' compared to Eq.~\eqref{eq:embcstokes4}.
\item
It is also possible to use the
Nitsche method, which instead of the positive term in
Eq.~\eqref{eq:embcstokes2}
it uses the following positive term:
\begin{equation}
   + \frac\alpha{h}(\vec v, \vec u)_{\Gamma_\text{D}} = \ldots
\label{eq:embcNitsche}
\end{equation}
where $h$ is a length scale representing the size of an element and
$\alpha$ is a factor that is supposed to be large enough. The Nitsche method is
easier to use (only a boundary integral, no volume integrals), but the optimal
factor is problem and mesh dependent ($h$).

\end{enumerate}
For the implementation we rewrite Eq.~\eqref{eq:embcstokes2} as
\begin{multline}
  \big(\ten D_v,2\eta\ten D\big)
    -(\nabla\cdot\vek v,p)
    -(\vek v,2\eta \ten D\cdot\vec n-p\vec n)_{\Gamma_{\text{D}}}%\\
 +\kappa M_{km}^{-1}(\vec v\vec n,\phi_k)_{\Gamma_\text{D}}:(\phi_m,\vek n\vec u
+\vec u\vec n)_{\Gamma_\text{D}}\\
 =\kappa M_{km}^{-1}(\vec v\vec n,\phi_k)_{\Gamma_\text{D}}:(\phi_m,\vek n\vec
u_\text{D}+\vec u_\text{D}\vec n)_{\Gamma_\text{D}}\\
+(\vek v,\vek t_{\text{N}})_{\Gamma_{\text{N}}}+(\vek v,\vek f)
\qquad\text{for all }\vek v
\label{eq:embcstokes3}
\end{multline}
Assuming $\phi$ is also the shape function for velocity, we write
$\vec v=\vec v_N\phi_N=v_N^i\phi_N\vec e_i$ and $\vec u=\vec
u_M\phi_M=u_M^j\phi_M\vec e_j$.
On element level we need to compute the element matrix $K_{NM}^{ij}$ from: 
\begin{equation}
   v^i_NK_{NM}^{ij}u^j_M=
     M^{-1}_{km}(\vec v\vec n,\phi_k)_{\Gamma_\text{D}}:(\phi_m,\vek n\vec u
        +\vec u\vec n)_{\Gamma_\text{D}}
\end{equation}
where $i,j=1,\dots,\texttt{ndim}$ and $N,M=1,\dots,\texttt{nnodes per element}$.
We find that
\begin{equation}
\begin{split}
   v^i_NK_{NM}^{ij}u^j_M&=
     \vec v_N\cdot\vec u_M M^{-1}_{km}(\phi_N\vec n,\phi_k)_{\Gamma_\text{D}}\cdot
(\phi_m,\vek n\phi_M)_{\Gamma_\text{D}}\\
&\quad+ \vec u_M\cdot(\phi_N\vec n,\phi_k)_{\Gamma_\text{D}}M^{-1}_{km}(\phi_m,\vek
n\phi_M)_{\Gamma_\text{D}}\cdot\vec v_N
\end{split}
\end{equation}
 Central in evaluating
$K_{NM}^{ij}$ is the matrix:
\begin{equation}
   S_{NM}^{ij} =
(\phi_Nn_i,\phi_k)_{\Gamma_\text{D}}M^{-1}_{km}(\phi_m,n_j\phi_M)_{\Gamma_\text{D}} 
\end{equation}
and thus
\begin{equation}
   K_{NM}^{ij}=S_{NM}^{mm}\delta_{ij}+S_{NM}^{ji}
\end{equation}
In the actual implementation we first compute:
\begin{equation}
   A_{mM}^i=(\phi_mn_i,\phi_M)
\end{equation}
and then
\begin{equation}
   B_{kM}^i=M^{-1}_{km}A_{mM}^i
\end{equation}
using Cholesky decomposition of $M_{km}$. The matrix $B_{kM}^i$ can also be
used in computing the right-hand side if $\vec u_\text{D}\neq\vec 0$.

\subsection{Weak rigid particle boundary conditions}
\label{sec:weakparticle}

(It is advised to first read Sec.~\ref{sec:weakDflow} on weak Dirichlet
boundary conditions for flow problems.

At the boundary $\Gamma_\text{p}$ of a freely floating rigid particle
the full Dirichlet condition cannot be
imposed, since the velocity still contains unknowns:
\begin{equation}
   \vek u = \vek U + \vek\omega \times (\vek x - \vek x_\text{p})
   \label{eq:velocitypar}
\end{equation}
where $\vek x_\text{p}$ is the center of mass of the particle,
$\vek U$ is the translational velocity of the center of mass and $\vek\omega$
the rotational rate of the particle. Both $\vek U$ and
 $\vek\omega$ are unknown and
need to be determined from the force and torque free condition.

The basic formulation is, as in Sec.~\ref{sec:weakDflow},
based on the open boundary formulation, where the open
boundary is applied to $\Gamma_\text{p}$, the particle boundary.
For example, for
Stokes flow we find from Eq.~\eqref{eq:obcstokes}:
\begin{equation}
  \big(\ten D_v,2\eta\ten D\big)
    -(\nabla\cdot\vek v,p)
    -(\vek v,2\eta \ten D\cdot\vec n-p\vec n)_{\Gamma_{\text{p}}}=
(\vek v,\vek t_{\text{N}})_{\Gamma_{\text{N}}}+(\vek v,\vek f)
\qquad\text{for all }\vek v
\label{eq:parbcstokes1}
\end{equation}
where $\vec n$ is the outwardly directed unit normal vector with respect to the
fluid (so $-\vek n$ is the outwardly directed unit normal vector on the
particle).

Introducing the variations $\vek V$ and $\vek\chi$ of $\vek U$ and
$\vek\omega$, respectively, we include the force and torque free conditions as
follows:
\begin{multline}
  \big(\ten D_v,2\eta\ten D\big)
    -(\nabla\cdot\vek v,p)
    -(\vek v-\vek V-\vek\chi\times(\vek x - \vek x_\text{p}),
          2\eta \ten D\cdot\vec n-p\vec n)_{\Gamma_{\text{p}}}=\\
(\vek v,\vek t_{\text{N}})_{\Gamma_{\text{N}}}+(\vek v,\vek f)
\qquad\text{for all }(\vek v, \vek V, \vek\chi)
\label{eq:parbcstokes2}
\end{multline}
For imposing the velocity condition on the particle 
Eq.~\eqref{eq:velocitypar} we add a positive definite term like in
Sec.~\ref{sec:weakDflow} with $\vek v$ replaced by $\vec v-\vek
V-\vek\chi\times(\vek x - \vek x_\text{p})$ and
$\vek u-\vek u_\text{D}$ by
$\vec u-\vek U-\vek\omega\times(\vek x - \vek x_\text{p})$:
\begin{multline}
  \big(\ten D_v,2\eta\ten D\big)
    -(\nabla\cdot\vek v,p)
    -(\vek v-\vek V-\vek\chi\times(\vek x - \vek x_\text{p}),
          2\eta \ten D\cdot\vec n-p\vec n)_{\Gamma_{\text{p}}}\\
\shoveleft{
 +\kappa M_{km}^{-1}((\vec v-\vek V-\vek\chi\times(\vek x - \vek x_\text{p}))
\vec n,\phi_k)_{\Gamma_{\text{p}}}:}\\
\shoveright{(\phi_m,[\vek n(\vec u-\vek U-\vek\omega\times(\vek x - \vek
x_\text{p}))+(\vec u-\vek U-\vek\omega\times(\vek x - \vek
x_\text{p}))\vec n])_{\Gamma_{\text{p}}}}\\
=(\vek v,\vek t_{\text{N}})_{\Gamma_{\text{N}}}+(\vek v,\vek f)
\qquad\text{for all }(\vek v, \vek V, \vek\chi)
\label{eq:parbcstokes3}
\end{multline}
Remarks:
\begin{enumerate}
   \item The method is consistent.
   \item For Stokes $\kappa=\eta$ is expected to be optimal.
   \item The original formulation \cite{choi_alignment_2012,Choi11b} uses
        an additional field $\ten E=\ten D(\vek u)$.
        In the formulation above the field $\ten E$
        has been eliminated to obtain the
        so-called primary formulation.
   \item Contrary to using Lagrange multipliers with additional unknowns
$(\vek U,\vek\omega)$, the diagonal is filled with positive values, helping
  iterative solvers.
   \item The full weak form is obtained by including the term (continuity eq.)
\begin{equation}
   \ldots -(q,\nabla\cdot\vek u) \dots
\end{equation}
to the left-hand side of the weak form.
   \item If a force $\vek F$ and $\vek M$ are imposed on the particle, the
the following terms must be added to the right-hand side
\begin{equation}
   \ldots = (\vek V,\vek F)+(\vek\chi,\vek M)\dots
\end{equation}
   \item The method is not limited to the Stokes problem and can be applied to
         other flow problems (Navier-Stokes, viscoelastic, \dots)
         by adding an identical additional term to the open
         boundary formulation for the problem at hand.
   \item The method is non-symmetric even for a symmetric problem, such as the
         Stokes problem. This can be remedied by adding the transposed open
         boundary term:
\begin{equation}
  \ldots  -(2\eta \ten D_v\cdot\vec n-q\vec n,
 \vek u-\vek U-\vek\omega\times(\vek x -
      \vek x_\text{p}))_{\Gamma_{\text{p}}} \dots
\end{equation}
to the left-hand side of the weak form.
   \item The positive (GLS) term is found by minimizing the positive
functional with respect to $\vek u$, $\vek U$ and $\vek\omega$:
\begin{multline}
 \frac14\kappa M_{km}^{-1}(\phi_k,[\vek n(\vec u-\vek U-\vek\omega\times(\vek x
- \vek
x_\text{p}))+(\vec u-\vek U-\vek\omega\times(\vek x - \vek
x_\text{p}))\vec n])_{\Gamma_{\text{p}}}:\\
(\phi_m,[\vek n(\vec u-\vek U-\vek\omega\times(\vek x - \vek
x_\text{p}))+(\vec u-\vek U-\vek\omega\times(\vek x - \vek
x_\text{p}))\vec n])_{\Gamma_{\text{p}}}\\
\end{multline}
   \item The positive term stabilizes the method. Another possibility for
stabilization is the Baumann-Oden approach, which adds a \emph{negative}
         transposed open
         boundary term:
\begin{equation}
  \ldots  +(2\eta \ten D_v\cdot\vec n-q\vec n,
 \vek u-\vek U-\vek\omega\times(\vek x -
      \vek x_\text{p}))_{\Gamma_{\text{p}}} \dots
\end{equation}
where the pressure part might be left out, since the method is non-symmetric
anyway.
  \item
It is also possible to use the
Nitsche method, which instead of the positive term in
Eq.~\eqref{eq:parbcstokes3}
it uses the following positive term:
\begin{equation}
   + \frac\alpha{h}(\vec v-\vek V-\vek\chi\times(\vek x - \vek x_\text{p}),
\vec u-\vek U-\vek\omega\times(\vek x - \vek
x_\text{p}))_{\Gamma_\text{p}} = \ldots
\label{eq:parbcNitsche}
\end{equation}
where $h$ is a length scale representing the size of an element and
$\alpha$ is a factor that is supposed to be large enough. The Nitsche method is
easier to use (only a boundary integral, no volume integrals), but the optimal
factor is problem and mesh dependent ($h$).
\item
In the implementation of Eq.~\eqref{eq:parbcstokes2} and
\eqref{eq:parbcstokes3} we can use that $\vek v$ and $\vek u$ are replaced by
$\vek v- \vek V-\vek\chi\times(\vek x-\vec x_\text{p})$ and
$\vek u- \vek U-\vek\chi\times(\vek x-\vec x_\text{p})$, respectively.
That is, originally we would have terms like
\begin{equation}
   (v_i,S_{ij}u_j)_{\Gamma_\text{p}}=v_i^k(\phi_k,S_{ij}\phi_m)_{\Gamma_\text{p}}u_j^m
\end{equation}
which now needs to be replaced by
\begin{equation}
   (v_i-V_i-N_{il}x_l,S_{ij}(u_j-U_j-L_{jn}x_n))_{\Gamma_\text{p}}=
 (v_i^k\phi_k-V_i-N_{il}x_l,S_{ij}(u_j^m\phi_m -U_j-L_{jn}x_n))_{\Gamma_\text{p}}
\label{eq:particlestable}
\end{equation}
where $L_{ij}$ is a coefficient matrix depending on the ``additional unknowns''
and $N_{ij}$ is the variation of $L_{ij}$. For example for a 3D particle we
have:
\begin{equation}
  \mat L=
  \begin{pmatrix}
    0 & -\omega_3 & \omega_2 \\ 
    \omega_3 & 0 & -\omega_1 \\ 
    -\omega_2 & \omega_1 & 0 
  \end{pmatrix},\quad
  \mat N=
  \begin{pmatrix}
    0 & -\chi_3 & \chi_2 \\ 
    \chi_3 & 0 & -\chi_1 \\ 
    -\chi_2 & \chi_1 & 0 
  \end{pmatrix}
\end{equation}
Once we have computed $(\phi_k,S_{ij}\phi_m)$, we can easily find the other
matrices
\[
\begin{split}
(1,S_{ij}\phi_m)_{\Gamma_\text{p}},\quad
(x_l,S_{ij}\phi_m)_{\Gamma_\text{p}}\quad
(\phi_k,S_{ij})_{\Gamma_\text{p}},\quad
(1,S_{ij})_{\Gamma_\text{p}},\\%\quad
(x_l,S_{ij})_{\Gamma_\text{p}},\quad
(\phi_k,S_{ij}x_n)_{\Gamma_\text{p}},\quad
(1,S_{ij}x_n)_{\Gamma_\text{p}},\quad
(x_l,S_{ij}x_n)_{\Gamma_\text{p}}
\end{split}
\]
by using the following properties of the shape functions $\phi_k(\vek x)$:
\begin{align}
   \sum_k\phi_k(\vek x)&=1,\quad\forall\vek x\quad\text{(partition of unity)}\\
   \sum_k\vek x^k\phi_k(\vek x)&=\vek x,\quad\forall\vek x\quad\text{(iso-parametric
element)}
\end{align}
For example,
\begin{align}
  (1,S_{ij}\phi_m)_{\Gamma_\text{p}}&=\sum_k(\phi_k,S_{ij}\phi_m)\\ 
  (x_l,S_{ij}\phi_m)_{\Gamma_\text{p}}&=\sum_k x_l^k(\phi_k,S_{ij}\phi_m)\\
  (1,S_{ij})_{\Gamma_\text{p}}&=\sum_m(1,S_{ij}\phi_m)\\ 
  (x_l,S_{ij}x_n)_{\Gamma_\text{p}}&=\sum_m (x_l,S_{ij}\phi_m)x^m_n
\end{align}
The contribution of the matrices with $\vek x$
the system matrix is determined by the
multiplication with $\mat L$ and $\mat N$, according to
Eq.~\eqref{eq:particlestable}.
\end{enumerate}

\subsection{Weak interface conditions for the Stokes equation}
\label{sec:weakinterfacestokes}

We consider the Stokes equation
in a region $\Omega=\Omega_1\cup\Omega_2$ with an internal interface
$\Gamma$ (see Figure~\ref{fig:interfaceregion1}):
\begin{alignat}{3}
 -&\nabla\cdot(2\eta \ten D)+\nabla p & &=\vek f,&&\qquad\text{in }\Omega\\
   &\nabla\cdot\vek u& &=0&&\qquad\text{in }\Omega
\end{alignat}
The boundary and interface conditions are
\begin{gather}
    \vec u = \vec u_D\text{ on }\Gamma_D\\
    -p\vec n+2\eta\ten D\cdot\vec n=\vec t_{\text{N}}
\text{ on }\Gamma_{\text{N}}\\
    \lj\vek u\rj = \vek 0\text{ on }\Gamma\label{eq:ucont}\\
    \lj-p\vec n+2\eta\ten D\cdot\vec n\rj=\vec
0\label{eq:fcont}
\end{gather}
where the boundary $\partial\Omega=\Gamma_D\cup\Gamma_{\text{N}}$ has
been split into a
Dirichlet and a Neumann part and
 $\lj\vek u\rj=\vek u_2-\vek u_1$, with $\vek u_i$ the value of $\vek u$ on the
side of $\Omega_i$.
\begin{figure}[htbp!]
   \begin{center}
   \includegraphics{interface_diffusion}
   \end{center}
   \caption{Domain $\Omega=\Omega_1\cup\Omega_2$ with an interface $\Gamma$.
The unit normal $\vek n$ has been defined on both the outside boundary
$\partial\Omega$ (where $\vek n$ is the outwardly directed unit normal) and
on the internal interface $\Gamma$ (where the direction of $\vek n$ is from
domain $\Omega_1$ towards $\Omega_2$).}
    \label{fig:interfaceregion1}
\end{figure}

Remarks:
\begin{enumerate}
   \item The interface conditions Eqs.~\eqref{eq:ucont} and \eqref{eq:fcont}
         specify continuity of
         the velocity $\vek u$ and traction vector across the
         interface $\Gamma$.
   \item The continuity of the normal velocity $\vek n\cdot\lj\vek u\rj=0$ 
         is related to the incompressibility of the fluid and the assumption
         that there is no storage of mass on the interface.
   \item It is possible to include non-zero jumps in the
         traction vector vector across $\Gamma$ (surface tension,
         Marangoni stresses), but this has not been done here
         for simplicity.
   \item It is possible to include non-zero jumps in the
         tangential velocity vector $\vek u-(\vek n\cdot\vek u)\vek n$
         vector across $\Gamma$ (slip layer), but this has not been done here
         for simplicity.
\end{enumerate}

The weak formulation is based on the open boundary formulation, where the open
boundary is applied to interface $\Gamma$, on both domains $\Omega_1$ and
$\Omega_2$:
\begin{align}
  \big(\ten D_v,2\eta\ten D\big)
    -(\nabla\cdot\vek v,p)
    -\big(\vek v_1,(2\eta \ten D\cdot\vec n-p\vec n)_1\big)_{\Gamma}&=
(\vek v,\vek t_{\text{N}})_{\Gamma_{\text{N}_1}}+(\vek v,\vek f)
\text{ on }\Omega_1
\label{eq:emicstokes1}\\
  \big(\ten D_v,2\eta\ten D\big)
    -(\nabla\cdot\vek v,p)
    +\big(\vek v_2,(2\eta \ten D\cdot\vec n-p\vec n)_2\big)_{\Gamma}&=
(\vek v,\vek t_{\text{N}})_{\Gamma_{\text{N}_2}}+(\vek v,\vek f)
\text{ on }\Omega_2
\label{eq:emicstokes2}
\end{align}
where 
the volume integrals are evaluated on $\Omega_1$ and 
$\Omega_2$, respectively. For stabilization of these equations
we use the approach as
described in \cite{Choi11} for non-Newtonian fluids. For this, we define the
(symmetric)
auxiliary tensor $\ten E=(\nabla\vek u+(\nabla\vek u)^T)/2$
on both sides of the interface:
\begin{align}
   (\ten H,\ten E)-(\ten H,\nabla\vek u) 
         + (\ten H\cdot\vek n,\vek u-\hat{\vek u}_1)_\Gamma =
       0\quad\text{on }\Omega_1 \label{eq:stabstokes1}\\
   (\ten H,\ten E)-(\ten H,\nabla\vek u) 
         - (\ten H\cdot\vek n,\vek u-\hat{\vek u}_2)_\Gamma =
       0\quad\text{on }\Omega_2 \label{eq:stabstokes2}
\end{align}
where $\hat{\vek u}_i$ are the values of $\vek u$
at the interface.
We also add additional flux terms to
 Eqs.~\eqref{eq:emicstokes1}-\eqref{eq:emicstokes2}:
\begin{align}
  \ldots -\kappa_1\big(\vek v_1,
       (2\ten E -\nabla \vek u-\nabla\vek u^T)_1\cdot\vek n\big)_{\Gamma}
              &=(v,f)\text{ on }\Omega_1
   \label{eq:stabstokes3}\\
  \ldots +\kappa_2\big(\vek v_2,
       (2\ten E -\nabla \vek u-\nabla\vek u^T)_2\cdot\vek n\big)_{\Gamma}
              &=(v,f)\text{ on }\Omega_2
   \label{eq:stabstokes4}
\end{align}
with two arbitrary constants $\kappa_1$ and $\kappa_2$.

In order to impose the interface conditions
Eqs.~\eqref{eq:ucont}-\eqref{eq:fcont} we introduce the following notation:
\begin{align}
   \{v\}_\omega=(1-\omega)v_1+\omega v_2,\quad \omega\in[0,1]
\end{align}
where $\omega$ is a weighting factor for the two sides of the interface.
Note, that $\{v\}_0=v_1$, $\{v\}_1=v_2$ and
$\{v\}_\frac12=\frac12(v_1+v_2)\equiv\{v\}$. The symmetric weighting $\{v\}$ is
commonly used in DG-methods for diffusion problems \cite{Arnold02}.

The method is defined by the specification of the numerical fluxes
$\hat{\vek u}_i$ and $(2\eta \ten D\cdot\vec n-p\vec n)_i\cdot\vek n$
at the interface.
According to the interface conditions Eqs.~\eqref{eq:ucont}-\eqref{eq:fcont} the
numerical fluxes need to be the same on both sides of the interface (flux
continuity). In
addition we assume the stabilizing fluxes 
$\kappa_i(2\ten E -\nabla \vek u-\nabla\vek u^T)_i\cdot\vek n$
to be continuous across the
interface as well. We define the following numerical fluxes on the interface:
\begin{align}
   \hat{\vek u}_1=\hat{\vek u}_2 &= \{\vek u\}_{1-\omega}\\
(2\eta \ten D\cdot\vec n-p\vec n)_i&\Rightarrow
   \{2\eta \ten D\cdot\vec n-p\vec n\}_\omega \\    
   \kappa_i(2\ten E -\nabla \vek u-\nabla\vek u^T)_i\cdot\vek n&\Rightarrow
            \{\kappa(2\ten E -\nabla \vek u-\nabla\vek u^T)\}_\omega\cdot\vek n
\end{align}
where $\Rightarrow$ means ``replace with''.\\
Remarks:
\begin{enumerate}
   \item For the symmetric case $\omega=\frac12$ we get:
\begin{align}
   \hat{\vek u}_1=\hat{\vek u}_2 &= \{\vek u\}\\
(2\eta \ten D\cdot\vec n-p\vec n)_i&\Rightarrow
   \{2\eta \ten D\cdot\vec n-p\vec n\}\\    
   \kappa_i(2\ten E -\nabla \vek u-\nabla\vek u^T)_i\cdot\vek n&\Rightarrow
            \{\kappa(2\ten E -\nabla \vek u-\nabla\vek u^T)\}\cdot\vek n
\end{align}
The fluxes at the interface are defined as the averaged fluxes
on both sides of the the interface.
   \item For the one-sided case $\omega=0$ (bias towards side 1), we get:
\begin{align}
   \hat{\vek u}_1=\hat{\vek u}_2 &= \vek u_2\\
(2\eta \ten D\cdot\vec n-p\vec n)_i&\Rightarrow
   (2\eta \ten D\cdot\vec n-p\vec n)_1\\    
   \kappa_i(2\ten E -\nabla \vek u-\nabla\vek u^T)_i\cdot\vek n&\Rightarrow
            (\kappa_1(2\ten E -\nabla \vek u-\nabla\vek u^T)_1\cdot\vek n
\end{align}
The value of $\vek u$ at the interface is now determined by the value at side 2,
whereas the traction is determined by the value at side 1.
\end{enumerate}

After substituting the numerical fluxes and eliminating the auxiliary variable
$\ten E$ (see \cite{Choi11}), we get the following weak form:
\begin{multline}
  \big(\nabla\vek v^T,\eta(\nabla\vek u+\nabla\vek u^T)\big)
    -(\nabla\cdot\vek v,p)
    -(q,\nabla\cdot\vek u)
    +\big(\lj\vek v\rj,\{\eta(\nabla\vek u+\nabla\vek u^T)
         \cdot\vec n-p\vec n\}_\omega\big)_{\Gamma}=\\
 +\kappa_1(1-\omega)^2 M_{km}^{(1)-1}\big(\lj\vec v\rj\vec
n,\phi^{(1)}_k\big)_{\Gamma}:\big(\phi^{(1)}_m,(\vek n\lj\vec u\rj
+\lj\vec u\rj\vec n)\big)_{\Gamma}\\
 +\kappa_2\omega^2 M_{km}^{(2)-1}\big(\lj\vec v\rj\vec
n,\phi^{(2)}_k\big)_{\Gamma}:\big(\phi^{(2)}_m,(\vek n\lj\vec u\rj
+\lj\vec u\rj\vec n)\big)_{\Gamma}\\
=(\vek v,\vek t_{\text{N}})_{\Gamma_{\text{N}}}+(\vek v,\vek f)
   \label{eq:weakstokesinterface}
\end{multline}
where $\phi^{(i)}_k$ is a shape function defined on all elements that have
contribution to $\Omega_i$ and
$M^{(i)}_{km}=(\phi^{(i)}_k,\phi^{(i)}_m)$.
We have used summation convection for dummy indices $k$ and $m$,
both running over all nodes.
The stabilization parameters $\kappa_1\geq0$ and $\kappa_2\geq0$
have a dimension
identical to a viscosity and can be chosen freely (at least one
must be non-zero), but a value
near the actual
viscosity is expected to gives the best results.

Remarks:
\begin{enumerate}
   \item The interpolation order of the shape function $\phi^{(i)}_k$
 needs to be at
least equal to the scalar interpolation and in the implementation we
 take identical shape functions.
   \item The inverse of the mass matrix $M^{(i)}_{km}$ needs to be computed.
 Since we
will use discontinuous shape functions $\phi^{(i)}_k$, this can be done at
element level.
   \item Although $\phi^{(i)}_k$ is defined in all elements,
 the contribution of the
additional term is only non-zero for elements that are crossed by the boundary.
   \item The method is consistent.
   \item For one-sided weighting $\omega=0$ or $\omega=1$, one of the
         stabilization terms is zero.
   \item At the time of writing it is unclear what is the best value for
         $\omega$. Some first results using fluids with different viscosity 
         show that the pressure is most sensitive
         and that best results are obtained if stabilization
         is biased towards the least viscous fluid:
         \[ \omega=
             \begin{cases}
                0 & \text{if }\eta_1<\eta_2 \\
                1 & \text{if }\eta_1>\eta_2
             \end{cases}
         \]
\item
The method is non-symmetric even for a symmetric problem.
This can be remedied
easily by adding the transposed boundary terms to the
weak momentum balance:
\begin{multline}
  \ldots
    +\big(\lj\vek v\rj,\{\eta(\nabla\vek u+\nabla\vek u^T)
         \cdot\vec n-p\vec n\}_\omega\big)_{\Gamma}
    +\big(\{\eta(\nabla\vek v+\nabla\vek v^T)
         \cdot\vec n-q\vec n,\lj\vek u\rj\}_\omega\big)_{\Gamma}\\
    \ldots=(\vek v,\vek t_{\text{N}})_{\Gamma_{\text{N}}}+(\vek v,\vek f)
   \label{eq:interfacestokestransposed}
\end{multline}
   \item The positive term stabilizes the method. Another possibility for
stabilization is the Baumann-Oden approach, which adds a \emph{negative}
         transposed open boundary term:
\begin{multline}
  \ldots
    +\big(\lj\vek v\rj,\{\eta(\nabla\vek u+\nabla\vek u^T)
         \cdot\vec n-p\vec n\}_\omega\big)_{\Gamma}
    -\big(\{\eta(\nabla\vek v+\nabla\vek v^T)
         \cdot\vec n-q\vec n,\lj\vek u\rj\}_\omega\big)_{\Gamma}\\
    \ldots=(\vek v,\vek t_{\text{N}})_{\Gamma_{\text{N}}}+(\vek v,\vek f)
   \label{eq:interfacestokesBaumann}
\end{multline}
The negative sign for the transposed pressure term is optional.
  \item
Stabilization can also be achieved by 
Nitsche's method \cite{Nitsche71}, which instead of the positive term
in \eqref{eq:embcdiffusion2} uses the following positive term:
\begin{equation}
   + \frac{K}{h}(\lj \vek v\rj,\lj \vek u\rj)_\Gamma = \ldots
   \label{eq:Nitschestokesinterface}
\end{equation}
where $h$ is a length scale representing the size of an element and
$K$ is a factor that is supposed to be large enough. Nitsche's method also
uses the symmetric variant in Eq.~\eqref{eq:interfacestokestransposed}.
\end{enumerate}

Some remarks on the implementation and use of the elements:
\begin{enumerate}
   \item The two sides of the interface are modelled using \xfem\ with the
     degrees of freedom on both sides defined as two layers (see
     Sec.~\ref{sec:layers}).
   \item The boundary flux term 
    $\big(\lj\vek v\rj,\{\eta(\nabla\vek u+\nabla\vek u^T)
         \cdot\vec n-p\vec n\}_\omega\big)_{\Gamma}$
in Eq.~\eqref{eq:weakstokesinterface} and
   the transposed of this term in Eqs.~\eqref{eq:interfacestokestransposed}
and \eqref{eq:interfacestokesBaumann}, are build in two steps to avoid
specifying the material properties on both sides of the interface at once.
For this, we write the flux term as:
\begin{multline}
    \big(\lj\vek v\rj,\{\eta(\nabla\vek u+\nabla\vek u^T)
         \cdot\vec n-p\vec n\}_\omega\big)_{\Gamma}=\\
    (1-\omega)\big(\vek v_2-\vek v_1,\eta(\nabla\vek u+\nabla\vek u^T)
         \cdot\vec n_1-p\vec n_1\big)_{\Gamma}\\
    +\omega\big(\vek v_1-\vek v_2,\eta(\nabla\vek u+\nabla\vek u^T)
         \cdot\vec n_2-p\vec n_2\big)_{\Gamma}
\end{multline}
where $\vek n_1=\vek n$, $\vek n_2=-\vek n$ (outside normal of $\Omega_1$,
$\Omega_2$, respectively). These terms can be build separately using the same
element, by defining the normal as the outside normal with respect to the
contributing side. Also, we need to specify which side (=layer) is being build.
The weighting factors $1-\omega$ and $\omega$ can be applied at the main
program level.
   \item The full element matrix of a positive stabilization term 
         (there are two in
         Eq.~\eqref{eq:weakstokesinterface} and only one in
          Eq.~\eqref{eq:Nitschestokesinterface}) are build in one call of
         \texttt{build\_system}.
\end{enumerate}

\subsection{Weak Dirichlet boundary conditions for convection-diffusion
            problems}
\label{sec:weakDdiffusion}

(It is advised to first read Sec.~\ref{sec:weakDflow} on weak Dirichlet
boundary conditions for flow problems.)

The basic formulation is based on the open boundary formulation, where the open
boundary is applied to $\Gamma_\text{D}$, the Dirichlet boundary. For example, for
the convection-diffusion-reaction equation we find from
Eq.~\eqref{eq:weakformdiffusion}:
\begin{equation}
   (v,\gamma\pderiv{c}t)+(v,\gamma \vek u\cdot\nabla c)
   +(\nabla v,\alpha\nabla c) + (v,\beta c)
             + (v,h_{\text{N}})_{\Gamma_{\text{N}}}
            -(v,\vek n\cdot[\alpha\nabla c])_{\Gamma_{\text{D}}}
              =(v,f)
   \label{eq:embcdiffusion1}
\end{equation}
where $\vec n$ is the outwardly directed unit normal vector and 
the volume integrals evaluated on the material body only. This does not lead us
anywhere, because the Dirichlet condition:
\begin{equation}
    c = \vec c_\text{D}\text{ on }\Gamma_\text{D}
\end{equation}
has not been imposed yet. We do that by adding to 
Eq.~\eqref{eq:embcdiffusion1} a symmetric positive term:
\begin{multline}
   (v,\gamma\pderiv{c}t)+(v,\gamma \vek u\cdot\nabla c)
   +(\nabla v,\alpha\nabla c) + (v,\beta c) +
             (v,h_{\text{N}})_{\Gamma_{\text{N}}}
            -(v,\vek n\cdot[\alpha\nabla c])_{\Gamma_{\text{D}}}\\
 +\kappa M_{km}^{-1}(v,\vec n\phi_k)_{\Gamma_\text{D}}\cdot(\phi_m\vek n, [c-
         c_\text{D}])_{\Gamma_\text{D}} =(v,f)
\label{eq:embcdiffusion2}
\end{multline}
where $\phi_k$ is a shape function defined on all elements and
$M_{km}=(\phi_k,\phi_m)$. 
We have used summation convection for dummy indices $k$ and $m$,
both running over all nodes.
The parameter $\kappa>0$ has a dimension
identical to a diffusion coefficient and can be chosen freely, but a value
near the actual
diffusion coefficient gives the best results.

Remarks:
\begin{enumerate}
   \item The interpolation order of the shape function $\phi_k$ needs to be at
least equal to the scalar interpolation and in the implementation we
 take identical shape functions.
   \item The inverse of the mass matrix $M_{km}$ needs to be computed. Since we
will use discontinuous shape functions $\phi_k$, this can be done at element
level.
   \item Although $\phi_k$ is defined in all elements, the contribution of the
additional term is only non-zero for elements that are crossed by the boundary.
   \item The method is consistent.
\item
The method is non-symmetric even for a symmetric problem.
This can be remedied
easily by adding the transposed boundary terms to the
weak momentum balance:
\begin{equation}
\ldots
            -(v,\vek n\cdot[\alpha\nabla c])_{\Gamma_{\text{D}}}
            -(\vek n\cdot[\alpha\nabla v],c- c_\text{D})_{\Gamma_{\text{D}}}
         \ldots =(v,f)
\label{eq:embcdiffusion3}
\end{equation}
which is only valid if $\alpha$ does not depend of $c$.
   \item The positive term stabilizes the method. Another possibility for
stabilization is the Baumann-Oden approach, which adds a \emph{negative}
         transposed open boundary term:
\begin{equation}
  \ldots
            -(v,\vek n\cdot[\alpha\nabla c])_{\Gamma_{\text{D}}}
            +(\vek n\cdot[\alpha\nabla v],c- c_\text{D})_{\Gamma_{\text{D}}}
  \ldots
         =(v,f)
\label{eq:embcdiffusion_bau}
\end{equation}
  \item
Stabilization can also be achieved by 
Nitsche's method \cite{Nitsche71}, which instead of the positive term
in \eqref{eq:embcdiffusion2} uses the following positive term:
\begin{equation}
   + \frac{K}{h}(v, c-c_\text{D})_{\Gamma_\text{D}} = \ldots
\end{equation}
where $h$ is a length scale representing the size of an element and
$K$ is a factor that is supposed to be large enough. Nitsche's method also
uses the symmetric variant in Eq.~\eqref{eq:embcdiffusion3}.

\end{enumerate}

\subsection{Weak interface conditions for diffusion problems}
\label{sec:weakinterfacediffusion}

We consider the diffusion-reaction\footnote{Convection is not yet
considered here. See Sec.~\ref{sec:weakinterfaceconvection}.}
equation in a region $\Omega=\Omega_1\cup\Omega_2$ with an internal interface
$\Gamma$ (see Figure~\ref{fig:interfaceregion}):
\begin{equation}
%    \gamma(\pderiv{c}t + \vek u\cdot\nabla c) -\nabla \cdot( \alpha \nabla c)
    \gamma\pderiv{c}t -\nabla \cdot( \alpha \nabla c)
+ \beta c = f\text{ on }\Omega
 \label{eq:interface_diffusion}
\end{equation}
where the coefficient $\alpha$ is either a scalar or a tensor. All the
coefficients ($\alpha$, $\beta$, $\gamma$ and $f$) may vary in
%coefficients ($\alpha$, $\beta$, $\gamma$, $\vek u$ and $f$) may vary in
space and time and may also depend on other quantities.
\begin{figure}[htbp!]
   \begin{center}
   \includegraphics{interface_diffusion}
   \end{center}
   \caption{Domain $\Omega=\Omega_1\cup\Omega_2$ with an interface $\Gamma$.
The unit normal $\vek n$ has been defined on both the outside boundary
$\partial\Omega$ (where $\vek n$ is the outwardly directed unit normal) and
on the internal interface $\Gamma$ (where the direction of $\vek n$ is from
domain $\Omega_1$ towards $\Omega_2$).}
    \label{fig:interfaceregion}
\end{figure}
The boundary and interface conditions\footnote{Initial conditions
need to be specified as well.} are
\begin{align}
    c &= c_D\text{ on }\Gamma_D\\
    -\vec n\cdot(\alpha\nabla c)&=h_{\text{N}}\text{ on }\Gamma_{\text{N}}\\
     \lj c \rj &= \bar i \text{ on }\Gamma\label{eq:ibar}\\
    -\vec n\cdot\lj\alpha\nabla c\rj&=\bar j\text{ on }\Gamma\label{eq:jbar}
\end{align}
where the boundary $\partial\Omega=\Gamma_D\cup\Gamma_{\text{N}}$ has
been split into a Dirichlet and a Neumann part and $\lj c\rj=c_2-c_1$ is the
jump of $c$ across the interface, with $c_i$ the value of $c$ on the side of
$\Omega_i$.

Remarks:
\begin{enumerate}
   \item The standard case is $\bar i=\bar j=0$. This specifies continuity of
         the primary variable $c$ and continuity of the fluxes across the
         interface $\Gamma$.
   \item The case $\bar j\neq0$ can be interpreted by writing
         Eq.~\eqref{eq:jbar} as follows:      
\[
    -\vec n_1\cdot(\alpha\nabla c)_1
    -\vec n_2\cdot(\alpha\nabla c)_2+\bar j=0\text{ on }\Gamma
\]
where $\vek n_1=\vek n$, $\vek n_2=-\vek n$ (outside normal of $\Omega_1$,
$\Omega_2$, respectively). This is the flux balance at the interface (sum of
fluxes towards the interface is zero).
 This means that $\bar j$ is a local flux
\emph{supplying} the component $c$ to the system at the
 interface $\Gamma$.
    \item The case $\bar i\neq0$ is more difficult to interpret. It represents
a high local gradient that is modeled as a jump. For example the mass or heat
 transfer across a thin membrane, although the jump should be coupled to the
flux.
    \item For a similar approach to embedded interfaces in 
          diffusion problems using \xfem\ combined with Nitsche's method, see
          \cite{dolbow2009,dolbow2010}.
\end{enumerate}

The weak formulation is based on the open boundary formulation, where the open
boundary is applied to interface $\Gamma$, on both domains $\Omega_1$ and
$\Omega_2$:
\begin{align}
   %(v,\gamma\pderiv{c}t)+(v,\gamma \vek u\cdot\nabla c)
   (v,\gamma\pderiv{c}t)
   +(\nabla v,\alpha\nabla c) + (v,\beta c)
             + (v,h_{\text{N}})_{\Gamma_{\text{N1}}}
            -\big(v_1,\vek n\cdot(\alpha\nabla c)_1\big)_{\Gamma}
              &=(v,f)\text{ on }\Omega_1
   \label{eq:opendiffusion1}\\
   %(v,\gamma\pderiv{c}t)+(v,\gamma \vek u\cdot\nabla c)
   (v,\gamma\pderiv{c}t)
   +(\nabla v,\alpha\nabla c) + (v,\beta c)
             + (v,h_{\text{N}})_{\Gamma_{\text{N2}}}
            +\big(v_2,\vek n\cdot(\alpha\nabla c)_2\big)_{\Gamma}
              &=(v,f)\text{ on }\Omega_2
   \label{eq:opendiffusion2}
\end{align}
where 
the volume integrals are evaluated on $\Omega_1$ and 
$\Omega_2$, respectively. For stabilization these equations
we use the approach as
described in \cite{Choi11} for non-Newtonian fluids. For this, we define the
auxiliary variable $\vek e=\nabla c$ on both sides of the interface:
\begin{align}
   (\vek h,\vek e-\nabla c) + (\vek n\cdot\vek h,c-\hat c_1)_\Gamma =
       0\quad\text{on }\Omega_1 \label{eq:stabdiffusion1}\\
   (\vek h,\vek e-\nabla c) - (\vek n\cdot\vek h,c-\hat c_2)_\Gamma =
       0\quad\text{on }\Omega_2 \label{eq:stabdiffusion2}
\end{align}
where $\hat c_1$ and $\hat c_2$ are the values of $c$ at the interface.
We also add additional flux terms to
 Eqs.~\eqref{eq:opendiffusion1}-\eqref{eq:opendiffusion2}:
\begin{align}
  \ldots -\kappa_1\big(v_1,\vek n\cdot(\vek e -\nabla c)_1\big)_{\Gamma}
              &=(v,f)\text{ on }\Omega_1
   \label{eq:stabdiffusion3}\\
   \dots +\kappa_2\big(v_2,\vek n\cdot(\vek e- \nabla c)_2\big)_{\Gamma}
              &=(v,f)\text{ on }\Omega_2
   \label{eq:stabdiffusion4}
\end{align}
with two arbitrary constants $\kappa_1$ and $\kappa_2$.

In order to impose the interface conditions
Eqs.~\eqref{eq:ibar}-\eqref{eq:jbar} we introduce the following notation:
\begin{align}
   \{v\}_\omega=(1-\omega)v_1+\omega v_2,\quad \omega\in[0,1]
\end{align}
where $\omega$ is a weighting factor for the two sides of the interface.
Note, that $\{v\}_0=v_1$, $\{v\}_1=v_2$ and
$\{v\}_\frac12=\frac12(v_1+v_2)\equiv\{v\}$. The symmetric weighting $\{v\}$ is
commonly used in DG-methods for diffusion problems \cite{Arnold02}.

The method is defined by the specification of the numerical fluxes\footnote{In
the DG literature the numerical fluxes are $\hat c_1$, $\hat c_2$, 
$(\alpha\nabla c)_1$ and $(\alpha\nabla c)_2$. However only continuity for
the normal component (if a vector or tensor) is required.}.
$\hat c_1$, $\hat c_2$, $\vek n\cdot(\alpha\nabla c)_1$,
$\vek n\cdot(\alpha\nabla c)_2$. 
According to the interface conditions Eqs.~\eqref{eq:ibar}-\eqref{eq:jbar} the
numerical fluxes need to be the same on both sides of the interface (flux
continuity), except for possible additional jumps $\bar i$ and $\bar j$. In
addition we assume the stabilizing fluxes 
$\kappa_1\vek n\cdot(\vek e- \nabla c)_1$
and $\kappa_2\vek n\cdot(\vek e- \nabla c)_2$ to be continuous across the
interface as well. We define the following numerical fluxes on the interface:
\begin{align}
   \hat c_1 &= \{c\}_{1-\omega}-(1-\omega)\bar i\\
   \hat c_2 &= \{c\}_{1-\omega}+\omega\bar i\\
   -\vek n\cdot(\alpha\nabla c)_1&\Rightarrow
            -\{\vek n\cdot(\alpha\nabla c)\}_\omega - \omega \bar j \\
   -\vek n\cdot(\alpha\nabla c)_2&\Rightarrow
            -\{\vek n\cdot(\alpha\nabla c)\}_\omega + (1-\omega)\bar j \\
   \kappa_1\vek n\cdot(\vek e- \nabla c)_1&\Rightarrow
            \{\kappa\vek n\cdot(\vek e- \nabla c)\}_\omega\\
   \kappa_2\vek n\cdot(\vek e- \nabla c)_2&\Rightarrow
            \{\kappa\vek n\cdot(\vek e- \nabla c)\}_\omega
\end{align}
where $\Rightarrow$ means ``replace with''.\\
Remarks:
\begin{enumerate}
   \item For the symmetric case $\omega=\frac12$ we get:
\begin{align}
   \hat c_1 &= \{c\}-\tfrac12\bar i\\
   \hat c_2 &= \{c\}+\tfrac12\bar i\\
   -\vek n\cdot(\alpha\nabla c)_1&\Rightarrow
            -\{\vek n\cdot(\alpha\nabla c)\} - \tfrac12 \bar j \\
   -\vek n\cdot(\alpha\nabla c)_2&\Rightarrow
            -\{\vek n\cdot(\alpha\nabla c)\} + \tfrac12\bar j \\
   \kappa_1\vek n\cdot(\vek e- \nabla c)_1&\Rightarrow
            \{\kappa\vek n\cdot(\vek e- \nabla c)\}\\
   \kappa_2\vek n\cdot(\vek e- \nabla c)_2&\Rightarrow
            \{\kappa\vek n\cdot(\vek e- \nabla c)\}
\end{align}
The fluxes are averaged across the interface and the ``jumps'' are equally
distributed on both sides of the interface.
   \item For the one-sided case $\omega=0$ (bias towards side 1), we get:
\begin{align}
   \hat c_1 &= c_2-\bar i\\
   \hat c_2 &= c_2\\
   -\vek n\cdot(\alpha\nabla c)_1&\Rightarrow
            -\vek n\cdot(\alpha\nabla c)_1 \\
   -\vek n\cdot(\alpha\nabla c)_2&\Rightarrow
            -\vek n\cdot(\alpha\nabla c)_1 + \bar j \\
   \kappa_1\vek n\cdot(\vek e- \nabla c)_1&\Rightarrow
            \kappa_1\vek n\cdot(\vek e- \nabla c)_1\\
   \kappa_2\vek n\cdot(\vek e- \nabla c)_2&\Rightarrow
            \kappa_1\vek n\cdot(\vek e- \nabla c)_1
\end{align}
The value of $c$ on the interface is now determined by the value at side 2,
whereas the flux is determined by the value at side 1.
\end{enumerate}

After substituting the numerical fluxes and eliminating the auxiliary variable
$\vek e$ (see \cite{Choi11}), we get the following weak form:
\begin{multline}
   (v,\gamma\pderiv{c}t)
   +(\nabla v,\alpha\nabla c) + (v,\beta c)
            +\big(\lj v\rj,\vek n\cdot\{\alpha\nabla c\}_\omega\big)_{\Gamma}\\
 +\kappa_1(1-\omega)^2 M_{km}^{(1)-1}(\lj v\rj,\vec
n\phi^{(1)}_k)_\Gamma\cdot(\phi^{(1)}_m\vek n,\lj c\rj-
         \bar i)_\Gamma \\
 +\kappa_2\omega^2 M_{km}^{(2)-1}(\lj v\rj,\vec
n\phi^{(2)}_k)_\Gamma\cdot(\phi^{(2)}_m\vek n,\lj c\rj-
         \bar i)_\Gamma\\
 =(v,f) - (v,h_{\text{N}})_{\Gamma_{\text{N}}}
                               +(\{v\}_{1-\omega},\bar j)_\Gamma
   \label{eq:weakdiffusioninterface}
\end{multline}
where $\phi^{(i)}_k$ is a shape function defined on all elements that have
contribution to $\Omega_i$ and
$M^{(i)}_{km}=(\phi^{(i)}_k,\phi^{(i)}_m)$.
We have used summation convection for dummy indices $k$ and $m$,
both running over all nodes.
The stabilization parameters $\kappa_1\geq0$ and $\kappa_2\geq0$
have a dimension
identical to a diffusion coefficient and can be chosen freely (at least one
must be non-zero), but a value
near the actual
diffusion coefficient is expected to gives the best results.

Remarks:
\begin{enumerate}
   \item The interpolation order of the shape function $\phi^{(i)}_k$
 needs to be at
least equal to the scalar interpolation and in the implementation we
 take identical shape functions.
   \item The inverse of the mass matrix $M^{(i)}_{km}$ needs to be computed.
 Since we
will use discontinuous shape functions $\phi^{(i)}_k$, this can be done at
element level.
   \item Although $\phi^{(i)}_k$ is defined in all elements,
 the contribution of the
additional term is only non-zero for elements that are crossed by the boundary.
   \item The method is consistent.
   \item For one-sided weighting $\omega=0$ or $\omega=1$, one of the
         stabilization terms is zero.
   \item At the time of writing it is unknown what influence
         non-symmetric weighting ($\omega\neq\frac12$) 
         and/or one of $\kappa_1$, $\kappa_2$ set to zero has on 
         accuracy and stability.
\item
The method is non-symmetric even for a symmetric problem.
This can be remedied
easily by adding the transposed boundary terms to the
weak momentum balance:
\begin{equation}
\ldots
        +\big(\lj v\rj,\vek n\cdot\{\alpha\nabla c\}_\omega\big)_{\Gamma}\\
    +\big(\vek n\cdot\{\alpha\nabla v\}_\omega,\lj u\rj-\bar i\big)_{\Gamma}\\
         \ldots =(v,f)
   \label{eq:interfacediffusiontranposed}
\end{equation}
which is only valid if $\alpha$ does not depend of $c$.
   \item The positive term stabilizes the method. Another possibility for
stabilization is the Baumann-Oden approach, which adds a \emph{negative}
         transposed open boundary term:
\begin{equation}
  \ldots
        +\big(\lj v\rj,\vek n\cdot\{\alpha\nabla c\}_\omega\big)_{\Gamma}\\
    -\big(\vek n\cdot\{\alpha\nabla v\}_\omega,\lj u\rj-\bar i\big)_{\Gamma}\\
  \ldots
         =(v,f)
   \label{eq:interfacediffusionBaumann}
\end{equation}
  \item
Stabilization can also be achieved by 
Nitsche's method \cite{Nitsche71}, which instead of the positive term
in \eqref{eq:embcdiffusion2} uses the following positive term:
\begin{equation}
   + \frac{K}{h}(\lj v\rj,\lj c\rj-\bar i)_\Gamma = \ldots
   \label{eq:Nitscheinterface}
\end{equation}
where $h$ is a length scale representing the size of an element and
$K$ is a factor that is supposed to be large enough. Nitsche's method also
uses the symmetric variant in Eq.~\eqref{eq:interfacediffusiontranposed}.
The method proposed in \cite{dolbow2009,dolbow2010} is $\omega=\tfrac12$
(symmetric weighting) and Nitsche's method\footnote{In
\cite{dolbow2009,dolbow2010} the sign of the boundary terms is opposite to
ours. Therefore, we think the weak form in \cite{dolbow2009,dolbow2010} is
in error.}.
\end{enumerate}

Some remarks on the implementation and use of the elements:
\begin{enumerate}
   \item The two sides of the interface are modelled using \xfem\ with the
     degrees of freedom on both sides defined as two layers (see
     Sec.~\ref{sec:layers}).
   \item The boundary flux term $\big(\lj v\rj,\vek n\cdot\{\alpha\nabla
c\}_\omega\big)_{\Gamma}$ in Eq.~\eqref{eq:weakdiffusioninterface} and
   the transposed of this term in Eqs.~\eqref{eq:interfacediffusiontranposed}
and \eqref{eq:interfacediffusionBaumann}, are build in two steps to avoid
specifying the material properties on both sides of the interface at once.
For this, we write the flux term as:
\begin{equation}
   \big(\lj v\rj,\vek n\cdot\{\alpha\nabla c\}_\omega\big)_{\Gamma}=
   (1-\omega)\big(v_2-v_1,\vek n_1\cdot(\alpha\nabla c)_1\big)_{\Gamma}
   +\omega\big(v_1-v_2,\vek n_2\cdot(\alpha\nabla c)_2\big)_{\Gamma}
\end{equation}
where $\vek n_1=\vek n$, $\vek n_2=-\vek n$ (outside normal of $\Omega_1$,
$\Omega_2$, respectively). These terms can be build separately using the same
element, by defining the normal as the outside normal with respect to the
contributing side. Also, we need to specify which side (=layer) is being build.
The weighting factors $1-\omega$ and $\omega$ can be applied at the main
program level.
   \item The full element matrix of a positive stabilization term 
         (there are two in
         Eq.~\eqref{eq:weakdiffusioninterface} and only one in
          Eq.~\eqref{eq:Nitscheinterface}) are build in one call of
         \texttt{build\_system}.
  \item In the implementation $\bar i$ is defined as the value of $c$ in
        layer 2 minus the value in layer 1. This means that if the transposed
flux term is used (symmetric version 
Eq.~\eqref{eq:interfacediffusiontranposed}
or Baumann-Oden Eq.~\eqref{eq:interfacediffusionBaumann}) the coefficient
representing $\bar i$ needs to be reversed when building the contribution to
the second side:
\begin{multline}
   \big(\vek n\cdot\{\alpha\nabla v\}_\omega,\lj u\rj-\bar i\big)_{\Gamma}=\\
  (1-\omega)\big(\vek n_1\cdot(\alpha\nabla v)_1, u_2-u_1-\bar i\big)_{\Gamma}
  +\omega\big(\vek n_2\cdot(\alpha\nabla v)_2,u_1-u_2+\bar i\big)_{\Gamma}
\end{multline}
An alternative is to use a negative sign for \texttt{factorvec} in the build.
For the positive stabilization terms in Eq.~\eqref{eq:weakdiffusioninterface}
and Eq.~\eqref{eq:Nitscheinterface} there is not a problem, since these are
build in ``one go''.
\end{enumerate}


\subsection{Weak interface boundary conditions for convection-diffusion
            problems}
\label{sec:weakinterfaceconvection}

We consider the convection-diffusion-reaction
equation in a region $\Omega=\Omega_1\cup\Omega_2$ with an internal interface
$\Gamma$ (see Figure~\ref{fig:interfaceregion}):
\begin{equation}
    \gamma(\pderiv{c}t + \vek u\cdot\nabla c) -\nabla \cdot( \alpha \nabla c)
+ \beta c = f\text{ on }\Omega
 \label{eq:interface_convection}
\end{equation}
where the coefficient $\alpha$ is either a scalar or a tensor. All the
coefficients ($\alpha$, $\beta$, $\gamma$, $\vek u$ and $f$) may vary in
space and time and may also depend on other quantities.
The boundary and interface conditions\footnote{Initial conditions
need to be specified as well.} are
\begin{align}
    c &= c_D\text{ on }\Gamma_D\\
    -\vec n\cdot(\alpha\nabla c)&=h_{\text{N}}\text{ on }\Gamma_{\text{N}}\\
     \lj c \rj &= 0 \text{ on }\Gamma\\
    -\vec n\cdot\lj\alpha\nabla c\rj&=\bar j\text{ on }\Gamma\label{eq:jbarc}
\end{align}
where the boundary $\partial\Omega=\Gamma_D\cup\Gamma_{\text{N}}$ has
been split into a Dirichlet and a Neumann part and $\lj c\rj=c_2-c_1$ is the
jump of $c$ across the interface, with $c_i$ the value of $c$ on the side of
$\Omega_i$.

Remarks:
\begin{enumerate}
   \item The difference with Sec.~\ref{sec:weakinterfacediffusion} is the
inclusion of a convection term $\gamma\vek u\cdot\nabla c$
and it is assumed that $c$ is continuous across the interface $\lj c\rj=0$.
Including $\lj c\rj\neq 0$ is possible, but this is more complicated, since
both jumps in molecular and convective fluxes are coupled at the interface.
   \item In this section we only consider the instationary and convection term
and refer to Sec.~\ref{sec:weakinterfacediffusion} for the other terms.
\end{enumerate}

We assume that $\gamma$ obeys a conservation equation:
\begin{equation}
    \pderiv{\gamma}t + \nabla \cdot(\vek u \gamma) = 0 \quad\text{in }\Omega
   \label{eq:gammaconvervation}
\end{equation}
For example, if $\gamma$ is the density (or proportionally to it), this
equation is fulfilled. Assuming the interface does not store or release
``mass'', this means that the normal component of $\gamma\vek u$ is continuous
 across the interface:
\begin{equation}
  \vek n\cdot \lj \gamma\vek u\rj = 0
  \label{eq:massfluxcontinuity}
\end{equation}
Remarks:
\begin{enumerate}
  \item If $\gamma$ is a constant, we have $\nabla\cdot\vek u=0$ in $\Omega$
 and $\vek n\cdot \lj\vek u\rj = 0$ on $\Gamma$. 
  \item If $\gamma$ is a continuous function across the interface, we have
$\vek n\cdot \lj\vek u\rj = 0$ on $\Gamma$.
\end{enumerate}

Using \eqref{eq:gammaconvervation} 
we can now write:
\begin{equation}
    \gamma(\pderiv{c}t + \vek u\cdot\nabla c) =
    \pderiv{(\gamma c)}t + \nabla \cdot(\gamma \vek u c)
   \label{eq:cconverservation}
\end{equation}
where $\gamma \vek u c$ is the \emph{convective flux}. Since 
$\vek n\cdot \lj \gamma\vek u\rj = 0$ and $\lj c\rj=0$, which is imposed by the
methods for diffusion in the Sec.~\ref{sec:weakinterfacediffusion}, the
convective flux is also continuous across the interface. Therefore, it doesn't
need to be imposed separately. However, we might choose to do it anyway using,
for example, the Lesaint-Raviart upwind techniques. For this, we partially
integrate the convective flux in the weak form:
\begin{align}
   (v,\pderiv{(\gamma c)}t)-(\nabla v,\gamma\vek u c)+
              (v,\gamma \vek n\cdot\vek u c)_\Gamma+\ldots
              &=(v,f)\text{ on }\Omega_1
   \label{eq:openconvection1}\\
   (v,\pderiv{(\gamma c)}t)-(\nabla v,\gamma\vek u c)-
              (v,\gamma \vek n\cdot\vek u c)_\Gamma+\ldots
              &=(v,f)\text{ on }\Omega_2
   \label{eq:openconvection2}
\end{align}
and demand (normal) convective flux continuity by
\begin{equation}
   \gamma\vek u c\Rightarrow \widehat{\gamma\vek u c}=
\begin{cases}
(\gamma \vek u c)_1&\text{ if }\vek u\cdot\vek n>0\\
(\gamma \vek u c)_2&\text{ if }\vek u\cdot\vek n<0
\end{cases}
\label{eq:upwinding}
\end{equation}
i.e.\ the convective flux is replaced by the ``upstream value''.

Remark:
\begin{itemize}
   \item We have assumed here that $\gamma>0$ and thus that
         due to Eq.~\eqref{eq:massfluxcontinuity}
         the sign of $\vek n\cdot\vek u$ has the same value on both sides of
         the interface.
\end{itemize}
Now we substitute $\widehat{\gamma\vek u c}$ for the boundary convective flux, 
use Eq.~\eqref{eq:cconverservation} again and find
\begin{align}
   (v,\gamma\pderiv{c}t + \gamma\vek u\cdot\nabla c)+
              (v_1,\widehat{\gamma \vek n\cdot\vek u c}
              -(\gamma \vek n\cdot\vek u c)_1)_\Gamma+\ldots
              &=(v,f)\text{ on }\Omega_1\\
   (v,\gamma\pderiv{c}t + \gamma\vek u\cdot\nabla c)-
              (v_2,\widehat{\gamma \vek n\cdot\vek u c}
              -(\gamma\vek n\cdot \vek u c)_2)_\Gamma+\ldots
              &=(v,f)\text{ on }\Omega_2
\end{align}
Now finally, we substitute the value for $\widehat{\gamma\vek u c}$ in
Eq.~\eqref{eq:upwinding} and find the weak form:
\begin{equation}
   (v,\gamma\pderiv{c}t + \gamma\vek u\cdot\nabla c)+
              (\hat v,\lj\gamma \vek n\cdot\vek u c\rj)_\Gamma
              +\ldots
              =(v,f)\text{ on }\Omega
   \label{eq:weakconvectionupwinding}
\end{equation}
where
\begin{equation}
   \hat v=
\begin{cases}
v_2&\text{ if }\vek u\cdot\vek n>0\\
v_1&\text{ if }\vek u\cdot\vek n<0
\end{cases}
\label{eq:testfunctionupwinding}
\end{equation}
i.e.\ the test function has been taken ``upstream'', or in other words, the
jump term contributes to the upstream interface side.

Remarks:
\begin{enumerate}
\item Since $\gamma\vek n\cdot\vek u$ is continuous, the jump term can be
simplified by $\lj\gamma \vek n\cdot\vek u c\rj=\gamma \vek n\cdot\vek u\lj
c\rj$.
\item In practice the inclusion of Lesaint-Raviart upwinding at the interface
doesn't seem to make much of a difference for diffusion-dominated problems,
although I did not do a convergence analysis yet.
The same is
true for convection-dominated problems and sufficiently refined meshes (mesh
P\'eclet $<1$). If not (mesh P\'eclet $>1$), Lesaint-Raviart upwinding might
make a difference, but boundary layers will not be resolved and stabilizing in
the domain is also necessary in order to get stable solutions away from the
interface. In the limit of no diffusion ($\alpha=0$), Lesaint-Raviart upwinding
is needed to couple the domains.
\end{enumerate}


\subsection{Flow around a periodic array of cylinders in a channel 
for the Stokes equation}
\label{sec:xfem_cylinder}

In this example we consider the flow around a periodic array of cylinders in a
channel (see Fig.~\ref{fig:periodic_array}). The height of the channel is 8 and
the length of the spatial periodic is also 8. The radius of the cylinder is 1.
The imposed flow rate is 10.
\begin{figure}[htbp!]
   \begin{center}
   \includegraphics[scale=0.4]{periodic_array}
   \end{center}
   \caption{Periodic array of cylinders confined between two walls}
    \label{fig:periodic_array}
\end{figure}
We solve the problem using four different methods:
\begin{enumerate}
  \item 
 The standard \fem\ method using a boundary-fitted mesh.
  \item
 A fictitious domain method using a regular mesh and a
weak constraint for the imposed zero velocity on the cylinder.
  \item
 The \xfem\ method also using a regular mesh and a
weak constraint for the imposed zero velocity on the cylinder.
  \item
 The \xfem\ method also using a regular mesh and a
weak Dirichlet boundary condition
for the imposed zero velocity on the cylinder.
\end{enumerate}

The example files are \texttt{body\_fitted1.f90}, \texttt{fic\_dom1.f90},
\texttt{xfem3.f90}
and \texttt{xfem8.f90}, respectively.
The example \texttt{xfem3.f90} 
uses the concept of a layer (see Sec.~\ref{sec:layers})
to define the degrees of freedom in the fluid.
The examples \texttt{xfem1.f90} and \texttt{xfem2.f90} are similar
to \texttt{xfem3.f90} but use Dirichlet conditions to remove the internal
degrees of freedom in the cylinder domain.
The difference between the examples \texttt{xfem1.f90} and \texttt{xfem2.f90}
is that the latter uses a more advanced integration scheme (see below).
It is recommended to use the concept of layers for implementation of \xfem\
(like example \texttt{xfem3.f90}) because of additional error checking and
the efficiency and ease of use for fluid interfaces.
The files can be obtained by typing at the
command line:
\begin{quote}
   \texttt{getexample xfem}
\end{quote}
The examples can be run by
\begin{verbatim}
     mesh_cylinder_in_box1 < meshinput1.txt
     body_fitted1
     fic_dom1
     xfem3
\end{verbatim}
The generated meshes are depicted in Fig.~\ref{fig:mesh_bf}.
\begin{figure}[htbp!]
   \begin{center}
   \includegraphics[scale=0.4]{mesh_bf}\hspace{1cm}
   \includegraphics[scale=0.4]{mesh_fd}
   \end{center}
   \caption{Boundary-fitted mesh (left) and regular mesh (80$\times$80) with an
object used for the boundary condition on the cylinder.}
    \label{fig:mesh_bf}
\end{figure}
In Figs~\ref{fig:pressureall} and \ref{fig:vorticityall} the results for the
pressure and vorticity are given. It is clear that the results for \xfem\ are
similar to the boundary-fitted \fem.
As expected the local solution results of the
fictitious domain method are very bad near the cylinder wall, since the standard
\fem\ approximation cannot describe local jumps in the field. 
The command 
\begin{quote}
   \texttt{gnuplot gnu1.gpi}
\end{quote}
will show plots of the velocity profile half way between the cylinders,
the values of the pressure and
vorticity on the cylinder wall for the three methods and the velocity on the
cylinder wall for the fictitious domain method and for $\xfem$.
It should be noted that the fictitious
domain method performs much better for more global quantities, such as the
drag computed using the Lagrange multipliers or the global motion of freely
floating particles. However, if accurate local field quantities near the jump
are important the fictitious domain cannot be used.
\begin{figure}[htbp!]
   \begin{center}
   \includegraphics[scale=0.4]{pressure_bf}\vspace{0.5cm}
   \includegraphics[scale=0.4]{pressure_fd}\vspace{0.5cm}
   \includegraphics[scale=0.4]{pressure_xf}
   \end{center}
   \caption{Pressure contours for the boundary-fitted, fictitious domain and
      \xfem\ approach}
    \label{fig:pressureall}
\end{figure}
\begin{figure}[htbp!]
   \begin{center}
   \includegraphics[scale=0.4]{vorticity_bf}\vspace{0.5cm}
   \includegraphics[scale=0.4]{vorticity_fd}\vspace{0.5cm}
   \includegraphics[scale=0.4]{vorticity_xf}
   \end{center}
   \caption{Vorticity contours for the boundary-fitted, fictitious domain and
      \xfem\ approach}
    \label{fig:vorticityall}
\end{figure}

Some general remarks:
\begin{itemize}
   \item A levelset function $\phi(\vec x)$ \cite{Osher03}
 has been used to determine whether points are
inside ($\phi<0$), outside ($\phi>0$) or on the
cylinder ($\phi=0$).
   \item
 A crucial part of \xfem\ is an accurate method for computing a function on a
part of the domain. We use a level set function to determine the integration
region together with the module eltree (see Sec.~\ref{sec:subdivide}),
which is used to divide the element into small quads.
The smallest quads near the
interface are less than 1\% of the size of element in the example.
\texttt{xfem1.f90}. Even so, this does not seem to
be an optimal and robust method to perform the integration:
\begin{enumerate}
   \item It is relatively expensive to generate the subdivision and a lot of
integration points are generated. Especially in 3D this does not seem to be
attractive.
   \item It is not very robust if only a small part of the element needs to be
integrated and very small quads are required. In the example \texttt{xfem1.f90} the regular mesh is slightly deformed to have a more optimal configuration.
\end{enumerate}
A better approach is to combine the eltree subdivision (quadtrees in 2D,
octrees in 3D) with a
division of the smallest quads into triangles (2D) and tetrahedrons (3D)
aligned with the integration boundary (\cite{Min07}). This has been implemented
in the module \texttt{submesh}.
The advantages of the eltree approach are that we do not have to
assume that the interface is straight within an element and we only need the
sign of the levelset function. For a further subdivision into
triangles/tetrahedron that do not
intersect the interface we need to have the
value of the levelset function also
to find the intersection of the interface with the boundary of an element.
The examples \texttt{xfem2.f90} and \texttt{xfem3.f90}
use the more advanced integration scheme using the combination of
quadtree approach (eltree) and the division into triangles (submesh).

\end{itemize}

Some remarks on the source code
\texttt{xfem1.f90}, \texttt{xfem2.f90} and \texttt{xfem3.f90}:
\begin{itemize}
   \item The levelset function is first set in all nodes of the mesh.
   \item Elements intersected by the interface are subdivided into small quads
and triangles using the add-on \texttt{subdivision}. 
\item The
decision of whether nodes and elements are included or excluded is based on the
value of the levelset function in the nodes.
This can lead to undetected interfaces in some elements.
An example is given in Figure~\ref{fig:wrong_interface}. 
Here the red part of
the interface is undetected based on the levelset value in the vertices of the
elements only. In the left element the levelset function value
has different signs and the left element is subdivided and the interface is
detected on subdivision level.
In the right element all levelset function values have the same sign
and the element is not subdivided and consequently the red part of the
interface is undetected. The interface is not closed (contains a hole) and the
interior and/or exterior region has a local flat face\footnote{Note, that if
no subdivision into quadtrees is used and only into triangles the problem does
not exist. It only appears if neighboring elements have a different level of
subdivision.}
\begin{figure}[htbp!]
   \begin{center}
   \includegraphics{wrong_interface}
   \caption{Undetected interface. Left element subdivided into 2$\times$2
subelements. Right element not subdivided.}
    \label{fig:wrong_interface}
   \end{center}
\end{figure}

This problem can be resolved by subdividing all elements within a short range of
the interface (given by the parameter \texttt{split\_threshold} in the code)
and the decision whether an element is crossed by an interface is postponed
until after the subdivision (see Figure~\ref{fig:correct_interface}.
Since this affects the structure of the matrix,
the subdivision must be done before the assembling takes place. 
The examples
\texttt{xfem5.f90} and \texttt{xfem6.f90} follow this approach and are
basically the same problems as \texttt{xfem2.f90} and \texttt{xfem3.f90},
respectively. Also the example \texttt{xfem8.f90} (weak Dirichlet) uses
this approach.
\begin{figure}[htbp!]
   \begin{center}
   \includegraphics{correct_interface}
   \caption{Detected interface. Left and right element subdivided into 2$\times$2
subelements.}
    \label{fig:correct_interface}
   \end{center}
\end{figure}
\item Subdivision using eltree is done on element level
while building the system using the user supplied routines
\texttt{set\_ninti\_user} and \texttt{set\_Gauss\_integration\_user} in
the module \texttt{stokes\_usergauss\_m}. An alternative could be
to build the eltrees for all elements before the build. The advantage would be
that it is possible to exclude elements and degrees of freedom
from the build and solve if, for example,
the integration area becomes too small.
Also for time-dependent problems, where the interface does not move, it is
faster to do the subdivision only once at the start of the simulation.
The examples
\texttt{xfem5.f90} and \texttt{xfem6.f90} follow this approach and are
basically the same problems as \texttt{xfem2.f90} and \texttt{xfem3.f90},
respectively. Also the example \texttt{xfem8.f90} (weak Dirichlet) uses
this approach.
   \item For plotting the \xfem\ data we first convert the computational mesh
    to a plotting mesh using \texttt{mesh\_convert} with a user subroutine
   \texttt{userelmesh} (see Sec.~\ref{sec:meshconvert}). The subdivision of
elements is again defined by the eltree subdivision, but the level of
subdivision is can be different from than in the building of the system to
avoid big meshes
for plotting. For interpolating the computational results onto the plotting
mesh we simply define an object that includes all plotting nodes and sample the
solution. Since this can be time consuming, we define an optimal number of
`blocks' (see Sec.~\ref{sec:blocks}) in the mesh to reduce the search time.

\end{itemize}

Some specific remarks on the source code \texttt{xfem3.f90}:
\begin{itemize}
   \item A layer is used to define the degrees of freedom of the fluid.
The layer is defined by a
\texttt{nodeset} including all nodes outside the cylinder and all nodes
of the elements crossed by the cylinder boundary.
   \item The system matrix is build by applying \texttt{build\_system} to
the layer only:
\begin{listing}{1}
  call build_system ( mesh, problem, sysmatrix, rhsd, &
    elemsub=stokes_elem, oldvectors=oldvectors, &
    coefficients=coefficients, buildvector=.false., layer=1 )
\end{listing}
\item
An additional build 
with zero valued element matrices is performed for elements not
fully within a layer (because of the additional connections generated there) to
completely fill the system matrix:
\begin{listing}{1}
  call build_system ( mesh, problem, sysmatrix, rhsd, &
    coefficients=coefficients, zeromatvec=.true., &
    addmatvec=.true., exclude_single_layer =.true. )
\end{listing}

\end{itemize}

Some specific remarks on the source code \texttt{xfem1.f90} and \texttt{xfem2.f90}:
\begin{itemize}
   \item Elements inside the cylinder that are excluded from assembling into
the system are simply build with zero integration points leading to a zero
element matrix and vector.
   \item For deleting nodes internal to the cylinder we define a
\texttt{nodeset} and a zero Dirichlet condition for these nodes (see
Sec.~\ref{sec:essnodeset}). This is allowed since these nodes are only
connected to elements that are not assembled in the system.
\end{itemize}
The approach used is \texttt{xfem1.f90} and \texttt{xfem2.f90} is not
recommended anymore and the approach using layers as in \texttt{xfem3.f90} is
preferred.

Some specific remarks on the source code \texttt{xfem8.f90}, which uses weakly
imposed Dirichlet boundary conditions: 
\begin{itemize}
   \item The eltree structures for all the elements near the interface are
supplied to the element routine using \texttt{oldvectors} (see
Sec.~\ref{sec:oldvectors}).
   \item For building the weak boundary conditions an
\texttt{elementset} (see Sec.~\ref{sec:elementsets}) is used.
\item
A small integration
area can be a problem and is related to numerical errors,
both from numerical integration and rounding errors.
Remarks:
         \begin{itemize}
            \item Degrees of freedom coupled to nodes of elements with small
                  integration area have a small stiffness and mass
                  and may lead to an ill-conditioned system matrix.
                  Also the mass matrix $M_{km}$ can become numerically
                  non-positive definite due to rounding errors. Using
                  integration rules having an order as high as
                  possible (preferably exact)
                  seems to help in shifting the problem to smaller critical
                  areas.
            \item The subdivision on small areas consists of triangles only and
                  requires a large number of points to be accurate.
                  For the computation of the stiffness matrix ($Q_2$
                  quadrilateral elements), terms like $\xi^2\eta^4$,
                  $\xi^3\eta^3$ and $\xi^4\eta^2$ need to be integrated.
                  This requires a sixth-order accurate scheme on a triangle
                  (12-point scheme) in order to be exact.
                  For the mass matrix $M_{km}$ even terms of $\xi^4\eta^4$
                  need to be integrated, requiring an eighth-order accurate
                  scheme (16-point scheme) to integrate this
                  exactly\footnote{For 3D problems there are terms like
                  $\xi^4\eta^4\zeta^4$,
                  requiring a twelfth-order accurate scheme (140-point
                  scheme).}.
                  Using triangular/tetrahedral elements instead of
                  quadrilateral/hexahedral elements would reduce the number of
                  points significantly, but this is not yet possible in \tfem.
                  Luckily, it is only necessary to use the high-order
                  integration if the actual integration area/volume is small.
                  For larger integration area/volume low-order rules can be
                  used.
         \end{itemize}
One possibility to circumvent small integration areas
is to fully discard
elements (and interface within that element)
from the system if the integration area is smaller than some
value (see Fig.~\ref{fig:remove_element}). This method is being employed in the
\texttt{xfem8} example.
Another more advanced, but also more involved,
method is a mesh-deformation technique as
described in \cite{Choi10} to maximize the integration areas.
A ``poor man's'' deformation method can be employed in the example.
\end{itemize}
\begin{figure}[htbp!]
   \begin{center}
   \includegraphics[scale=1.0]{remove_element}
   \end{center}
   \caption{If the integration area becomes small, the stiffness and mass
matrix contributions to the virtual degrees in the red node can become very
small compared to the other nodes. This leads to pronounces
rounding errors in solving
the system and failure of the inversion of the mass matrix $M_{km}$ due to
becoming non-positive. Using exact integration rules helps to shift the problem
to smaller areas. A method for solving the problem for very small areas is
discarding 
the virtual/extended
degrees of freedom of the element (in the red node) and the contribution to the
system matrix of the element. Also the interface contribution to
the weak Dirichlet boundary condition of the element (red interface part)
is discarded.}
    \label{fig:remove_element}
\end{figure}

Examples \texttt{xfem9.f90} and \texttt{xfem10.f90} are derived from
\texttt{xfem8.f90} with the following differences:
\begin{itemize}
   \item \texttt{xfem9.f90} uses the symmetric variant given by 
Eq.~\eqref{eq:embcstokes4}. 
Optionally it is possible to use the
Nitsche method Eq.~\eqref{eq:embcNitsche}.
\item \texttt{xfem10.f90} uses the Baumann-Oden approach
Eq.~\eqref{eq:embcstokes6}.
\end{itemize}

Example \texttt{xfem12.f90} is an extension to 3D of
a combination of \texttt{xfem8.f90} and \texttt{xfem9.f90}.

\subsection{Flow around a freely floating particle in a channel 
for the Stokes equation}
\label{sec:xfem_particle}

This is example is similar to the flow around cylinder in
Sec.~\ref{sec:xfem_cylinder}, but now the fixed cylinder is replaced with a
freely floating cylinder. It is implemented using a weak embedded particle
condition as explained in Sec.~\ref{sec:weakparticle}. 

Examples \texttt{xfem13.f90} (2D circular particle) and \texttt{xfem14.f90} (3D
spherical particle) are derived from
\texttt{xfem8.f90+xfem9.f90} and \texttt{xfem12.f90}, respectively.
Example is like \texttt{xfem13.f90}, but now with an elliptical particle.
Note, that
only the solutions at start-up are computed,
since the particles are not tracked in time.

\subsection{Flow around a periodic array of cylindrical/spherical shaped drops
in a channel for the Stokes equation. Interface constraint.}
\label{sec:drop1}

This example is derived from the example in Sec.~\ref{sec:xfem_cylinder}.
However, now the cylinder is replaced by a fluid with a different viscosity
(10$\times$ as high as the surrounding fluid) and is put off-center.
So we consider the flow around a periodic array of drops in a
channel. The height of the channel is 8 and
the length of the spatial periodic is also 8. The radius of the drop is 1.
The imposed flow rate is 10. This is basically a time-dependent problem, since the interface
between the two fluids will evolve in time. 
We will not follow the interface here and thus basically solve only the 
initial velocity/pressure field at the start of this process.

We use the \xfem\ method as described in Sec.~\ref{sec:xfem_interpretation}.
The inside and outside fluid is put into separate layers of degrees of freedom.
The velocity degrees of freedom in the two layers are coupled at the
interface:
\begin{equation}
   \vek u_\text{inside} = \vek u_\text{outside} 
   \label{eq:velocitycont}
\end{equation}
The tractions on the interface are also the same on both sides of the
interface (except for a change in sign if the normal vectors taken been taken
opposite).
The pressures are not coupled and are allowed to jump at the interface.

The example uses a Lagrange multiplier to impose the continuity of
velocity Eq.~\eqref{eq:velocitycont} as a weak constraint. The Lagrange
multiplier is the traction at the interface.
The example file is \texttt{xfem4.f90}.
Like the
\texttt{xfem5.f90} and \texttt{xfem6.f90} examples 
an example \texttt{xfem7.f90} exists that
follows the approach of building the eltree for elements near the interface
before the definition and actual build of the system.

Some remarks on the source code (most of remarks in Sec.~\ref{sec:xfem_cylinder}
also apply here):
\begin{itemize}
   \item Two layers are defined.
A layer is defined by a
\texttt{nodeset} including all nodes
of the elements crossed by the cylinder boundary and either
 all nodes outside the cylinder (outside fluid) or 
all nodes inside the cylinder (inside fluid).
A global variable (\texttt{intregion}) is used to signal the
integration routines whether the integration region is inside or outside.
   \item The system matrix is build by applying \texttt{build\_system} separately
to the two layers.
\item
An additional build 
with zero valued element matrices is performed for elements not
fully within a single layer only
(because of the additional connections generated there) to
completely fill the system matrix. This includes all
elements at the interface and the `buffer' elements
between the interface elements and the single layer elements.
\item 
It is very tricky to find a suitable mesh and some optimization is necessary
(see \cite{Choi10}). In the example a suitable mesh is found by trial and error (by moving the center
of the cylinder and deforming the mesh).
Note, that starting with the example \texttt{xfem8.f90} mesh optimization
is not necessary anymore, since small integration area are discarded.
\item
The definition of the constraint for coupling the velocities couples two
objects, but both objects refer to the same object. The layers are different however:
\begin{listing}{1}
! define constraint on the object for coupling velocities between the layers

  call define_constraint ( mesh, input_probdef, object=1, layer1=1, object2=1,&
    layer2=2, physq=1, discretization='weak', elementdof=[2,0,2] )
\end{listing}

\end{itemize}

After running the code the results in Figure~\ref{fig:xfem4} can be obtained.
\begin{figure}[htbp!]
   \begin{center}
   \includegraphics[scale=0.4]{pressure_xf4}\vspace{0.5cm}
   \includegraphics[scale=0.4]{velocity_y_xf4}\vspace{0.5cm}
   \includegraphics[scale=0.4]{vorticity_xf4}
   \end{center}
   \caption{Pressure, vertical velocity and vorticity contours
      for a cylindrical drop using the \xfem\ approach.}
    \label{fig:xfem4}
\end{figure}
Notice the continuity of the velocity and the discontinuity of
the pressure and vorticity at the interface between the two fluids.
There is also a slight asymmetry in the contour lines of the pressure, which is
related to the slight off-center position of the drop in $x$-direction.
Notice also that the drop is rotating (look at the sign of the
velocity in vertical direction).

\subsection{Flow around a periodic array of cylindrical/spherical shaped drops
in a channel for the Stokes equation. Weak (embedded) interface condition.}
\label{sec:drop2}

The same problem as in Sec.~\ref{sec:drop1}
has been solved in example \texttt{xfem16.f90} using a weak interface
condition (see Sec.~\ref{sec:weakinterfacestokes}).
The problem has also been extended to 3D in example
\texttt{xfem17.f90}.

Remarks:
\begin{itemize}
   \item Both the unsymmetric version Eq.~\eqref{eq:weakstokesinterface}
 and symmetric version Eq.~\eqref{eq:interfacestokestransposed} can be used.
\item
Optionally it is also possible to use the
Baumann-Oden method Eq.~\eqref{eq:interfacestokesBaumann}.
\item
Optionally it is also possible to use the
Nitsche method Eq.~\eqref{eq:Nitschestokesinterface}.
\end{itemize}

\subsection{Flow around a periodic array of cylinders in a channel 
for a generalized Newtonian fluid}
\label{sec:xfem_cylinder_GN}
 
(This example has been developed by Arash Sarhangi Fard of the TU/e.)

The problem is in source file \texttt{xfem11.f90} and
is similar to the \texttt{xfem2.f90} problem
(see Sec.~\ref{sec:xfem_cylinder}), but now the fluid is replaced by a
generalized Newtonian fluid (Carreau fluid).

\subsection{A steady diffusion problem on a unit square/cube
            with a circular/spherical hole
            with Dirichlet
            and Neumann boundary conditions. Constant or varying diffusion
            coefficient. Weak Dirichlet boundary conditions.}

\label{sec:diffusion5}

We consider the steady
diffusion problem
on a unit square with a circular hole. The example is derived from
the \texttt{diffusion1} problem in Sec.~\ref{sec:diffusion1}.
The unit domain is cut by a circle and on the circle the scalar is imposed
using a weak Dirichlet boundary condition. The scalar is solved using \xfem\
with a weak embedded Dirichlet condition (see Sec.~\ref{sec:weakDdiffusion})

The source 
in the file \texttt{diffusion5.f90}
 and can be obtained by typing at the
command line:
\begin{quote}
   \texttt{getexample diffusion}
\end{quote}

The xfem method resembles the method in the flow problems
\texttt{xfem8.f90} and \texttt{xfem9.f90} in
Sec.~\ref{sec:xfem_cylinder}. 

An extension of the problem to 3D (unit cube, spherical hole)
is in the file \texttt{diffusion6.f90}.

Remarks:
\begin{itemize}
   \item Both the unsymmetric version Eq.~\eqref{eq:embcdiffusion2}
 and symmetric version Eq.~\eqref{eq:embcdiffusion2} can be used.
\item
Optionally it is also possible to use the
Nitsche method \cite{Nitsche71}.
\end{itemize}

\subsection{A steady diffusion problem on a unit square/cube
            with a circular/spherical region with a different diffusion
            constant. Weak interface conditions.}

\label{sec:diffusion7}

We consider the steady
diffusion problem
on a unit square with a circular region having different properties.
The interface conditions are imposed weakly.
The scalar is solved using \xfem\ with a weak embedded interface condition
(see Sec.~\ref{sec:weakinterfacediffusion}).

The source 
in the file \texttt{diffusion7.f90}
 and can be obtained by typing at the
command line:
\begin{quote}
   \texttt{getexample diffusion}
\end{quote}

Some results are shown in Figure~\ref{fig:diffusion7}.
\begin{figure}[htbp!]
   \begin{center}
   \includegraphics[scale=0.4]{scalar_df7}\vspace{0.5cm}
   \includegraphics[scale=0.4]{gradc_x}\vspace{0.5cm}
   \includegraphics[scale=0.4]{gradc_y}
   \end{center}
   \caption{Scalar $c$, $\plderiv{c}{x}$ and $\plderiv{c}{y}$ 
      contours for the \texttt{diffusion7} problem. Value of $\alpha=1$ in the
      outer region and $\alpha=0.05$ in the inner (circular) region. Mesh is
        $50\times50$ $Q_2$ elements.}
    \label{fig:diffusion7}
\end{figure}

An extension of the problem to 3D (unit cube, spherical region)
is in the file \texttt{diffusion8.f90}.

Remarks:
\begin{itemize}
   \item Both the unsymmetric version Eq.~\eqref{eq:weakdiffusioninterface}
 and symmetric version Eq.~\eqref{eq:interfacediffusiontranposed} can be used.
\item
Optionally it is also possible to use
Nitsche's method \cite{Nitsche71} and the Baumann-Oden approach.
\end{itemize}

\subsection{A steady convection-diffusion problem on a unit square
            with a circular region. At the interface ``heat'' is supplied.
            Weak interface conditions.}

\label{sec:convectiondiffusion8}

We consider the steady
convection-diffusion problem
on a unit square with a circular region. ``Heat'' is supplied at the circular
interface between the regions.
The interface conditions are imposed weakly 
(see Sec.~\ref{sec:weakinterfaceconvection}).
The scalar is solved using \xfem.

The source 
in the file \texttt{convection\_diffusion8.f90}
 and can be obtained by typing at the
command line:
\begin{quote}
   \texttt{getexample convection\_diffusion}
\end{quote}

Some results are shown in Figure~\ref{fig:convection_diffusion8}.
\begin{figure}[htbp!]
   \begin{center}
   \includegraphics[scale=0.4]{scalar_cd8}\vspace{0.5cm}
   \end{center}
   \caption{Scalar $c$ 
      contours for the \texttt{convection\_diffusion8} problem. The velocity
vector $\vek u=(10,5)$ is constant throughout the domain. At the circular
      interface $\bar{j}=100$.
      Mesh is $50\times50$ $Q_2$ elements.}
    \label{fig:convection_diffusion8}
\end{figure}

Remarks:
\begin{itemize}
   \item Lesaint-Raviart upwinding is optional.
   \item Both the unsymmetric version Eq.~\eqref{eq:weakdiffusioninterface}
 and symmetric version Eq.~\eqref{eq:interfacediffusiontranposed} can be used.
\item
Optionally it is also possible to use
Nitsche's method \cite{Nitsche71} and the Baumann-Oden approach.
\end{itemize}

\subsection{Extrudate swell of a Newtonian fluid}
\label{sec:swellNewtonian_xfem}

(This example has been developed by Young Joon Choi of the TU/e.)

The example is derived from a similar example using an ALE approach in
Sec.~\ref{sec:swellNewtonian}.
The fluid is Newtonian.
Use is made of the add-on surface advection (Sec.~\ref{sec:surfaceadvection}).
See \cite{Choi11c,Choi11b} and the references therein for the details on
solving the surface advection equation.

The example file is \texttt{extrudate\_swell4.f90} and can
be obtained by typing at the
command line:
\begin{quote}
   \texttt{getexample extrudate\_swell}
\end{quote}
The example uses an open boundary condition at the inflow boundary.
We assume zero normal traction and zero vertical velocity at the exit.
The free surface is traction free.
The flow is generated by imposing a flow rate constraint on the inflow
boundary.
The flow can be chosen to be planar of axisymmetric.

You can run the example by
\begin{quote}
   \texttt{mesh\_extrudate\_swell\_xfem < meshinput\_xfem.txt}\\
   \texttt{extrudate\_swell4 < input\_extrudate\_swell4.txt}\\
   \texttt{post\_swell4}
\end{quote}
Inspect the source files to see which data and plotfiles are generated.
Results of the final free surface for a planar flow are shown in
Fig.~\ref{fig:swell_newtonian_xfem}.
\begin{figure}[htbp!]
   \begin{center}
   \includegraphics[scale=0.8]{mesh_swell_0300}
   \end{center}
   \caption{Final surface 
    of the extrudate swell problem using a Newtonian fluid.}
    \label{fig:swell_newtonian_xfem}
\end{figure}

In the example there is an option to add surface tension. The implementation is
based on a membrane with an isotropic surface stress $\gamma$. This means,
we add to the right-hand side of the weak form:
\begin{equation}
  -\int_S(\nabla\vek v)^T:\gamma(\ten I-\vek n\vek n)\D S+
\int_{\partial S}\vek v\cdot\gamma\vek n_\text{c}\D s
\end{equation}
where $\vek n$ is the unit normal vector on the surface $S$, $\partial S$ the
boundary of the surface (a point in 2D and a curve in 3D) and
$\vek n_\text{c}$ the
unit vector that is tangential to $S$ and normal to $\partial S$ and pointing
``outwards''.


\subsection{Flow around a cylinder confined between two walls}
\label{sec:cylinderve_xfem}

(This example has been developed by Young Joon Choi of the TU/e.)

This problem is the 
flow of a viscoelastic fluid (Oldroyd-B) around a cylinder confined between
two walls.
The entry and exit sides are connected by periodical boundary conditions
and a constant flow rate is imposed.
The cylinder is modelled using \xfem\ and weak boundary conditions.
The problem is similar to the one described in \cite{Choi11}.

The example file is \texttt{cylinder1.f90} and
and can be obtained by typing at the
command line:
\begin{quote}
   \texttt{getexample confined\_cylinder\_xfem}
\end{quote}
You can run the example by
\begin{quote}
   \texttt{mesh\_channel\_three\_parts < meshinput.txt}\\
   \texttt{cylinder1 < input\_extrudate\_swell1.txt}\\
   \texttt{post\_cylinder1}
\end{quote}
Inspect the source files to see which data and plotfiles are generated.
